\documentclass[a4paper, 10pt]{article}
\usepackage[utf8]{inputenc}
\usepackage[italian]{babel}
\usepackage[margin=0.8in]{geometry}
\usepackage{graphicx}
\usepackage{float}
\usepackage[table,xcdraw]{xcolor}
\usepackage[hidelinks]{hyperref}

\begin{document}
\begin{center}
\includegraphics[width=0.3\textwidth]{../template/images/logo.png}
\hspace{3cm}
\includegraphics[width=0.2\textwidth]{../template/images/logoUnipd.png}\\
\end{center}
\begin{flushright}
    \
    \textbf{}\\
    Padova - XX Maggio 2022
\end{flushright}
Egregio Prof. Vardanega,\\
Egregio Prof. Cardin,\\\\
Con la presente il gruppo Yakuzaishi intende comunicarvi la candidatura alla revisione relativa la Product Baseline, al fine di esporre l'avanzamento dello sviluppo del progetto denominato:
\begin{center}
    \textbf{ShopChain}
\end{center}
proposto dall'azienda Sync Lab S.r.l.\\\\
Nel repository all'indirizzo:
\begin{center}
    \href{https://yakuzaishi-swe.github.io/docs/PB/}{\underline{https://yakuzaishi-swe.github.io/docs/PB/}}\\ 
\end{center}
potrà trovare tutti i documenti che trattano in maniera approfondita gli aspetti di pianificazione, normazione, qualifica e analisi.\\\\
Nella cartella \textit{Documenti Esterni} saranno presenti i file:
\begin{itemize}
    \item \textit{analisi\_dei\_requisiti\_v3.0.0};
    \item \textit{piano\_di\_progetto\_v2.0.0};
    \item \textit{piano\_di\_qualifica\_v2.0.0};
    \item \textit{glossario\_v2.0.0};
    \item \textit{manuale\_utente\_v1.0.0};
    \item \textit{allegato\_tecnico\_v1.0.0}.
\end{itemize}
Nella cartella \textit{Documenti Interni} saranno presenti i file:
\begin{itemize}
    \item \textit{norme\_di\_progetto\_v2.0.0}.
\end{itemize}
Saranno inoltre presenti, nelle rispettive cartelle, i seguenti verbali di periodo:
\begin{itemize}
    \item \textit{verbale\_esterno\_2022\_03\_28};
    \item \textit{verbale\_esterno\_2022\_05\_02};
    \item \textit{verbale\_interno\_2022\_05\_03}.
\end{itemize}
Per completezza, inseriamo di seguito la lista dei membri del gruppo.
\begin{table}[H]
    \centering
    \renewcommand{\arraystretch}{1.8}
    \rowcolors{2}{green!100!black!40}{green!100!black!30}
    \begin{tabular}{l|cccccc|c}
      \rowcolor[HTML]{125E28} 
      \multicolumn{1}{c}{\color[HTML]{FFFFFF}\textbf{Nome}} 
      & \color[HTML]{FFFFFF}\textbf{Matricola}\\
      \hline
      Bugno Francesco & 1170991\\
      Busacca Luca & 1227589\\
      Carturan Luca & 1094033\\
      Filosofo Michele & 1094325\\
      Furlan Dario & 1223861\\
      Mattarello Francesco & 1172722\\
      Midena Matteo & 1227272\\
      \hline             
    \end{tabular}
  \end{table}
\textbf{}\\
Cordiali saluti,\\
\begin{flushright}
    \
    Responsabile di Progetto\\
    Francesco Bugno \hspace{1.1cm}\hspace{1cm}
\end{flushright}
\end{document}
